% Chapter 1 — The Objective and the Commitments (KLEPPMANN, 4 pp)
% Discharges constitution §1, §2.

\begin{abstract}
Reconciliation exists because many systems store the same facts, and independently
maintained copies eventually disagree. The Ledger removes the cause: every fact is
recorded exactly once, in an immutable log of atomic transactions, and every other
number --- balances, valuations, profit and loss, reports --- is a deterministic
projection of that single record, reproducible by any party and reconstructible at any
past date. This specification defines the system in full, rated one-way against the
Ledger Framework Constitution v1.1: the objects the system manipulates and the signed
coordinate vector every balance generalises to; the three machines that evaluate the
pipeline --- the Event Monitor, the Events Executor, and the Transaction Executor,
three roles and one writer; smart contracts as pure, stateless transcriptions of legal
terms; the three homes of state; valuation and profit and loss; ingestion, corporate
actions, and the market data operator; collateral, margin, and lending under the
declared legal regime of the governing agreement; virtual ledgers and the total return
swap; the settlement interface; presentation and reporting projections; alignment with
the Common Domain Model; and the invariant catalogue and executable-check regime that
make every guarantee a property a machine can run. Six instruments with fixed terms
and fixed quantities recur through every chapter, so that each claim is exercised on
the same worked facts it governs.
\end{abstract}

\tableofcontents
\clearpage

\section{The Objective and the Commitments}\label{ch:objective}

This chapter discharges \S1 and \S2 of the constitution.

\subsection{One record, or reconciliation}

The objective of the Ledger is a clean post-trade and risk-management platform: one
system of record from which positions, valuations, profit and loss, collateral, and
reports are all derived --- exactly, deterministically, and reproducibly.

The obstacle is the architecture of the infrastructure the industry runs today.
Post-trade processing is not one system but many --- trading systems, risk engines,
general ledgers, custodians, regulatory engines --- and each holds its own copy of
positions, cash balances, corporate-action adjustments, and lifecycle states, updated
on its own schedule by its own logic. The copies diverge, and reconciliation ---
detecting and repairing the divergences --- is one of the largest operational costs
and risks in the industry. The root cause is architectural: \textbf{the same fact is
stored more than once.} Two stores of one fact maintained by independent logic will
eventually disagree, and no amount of process removes a failure mode the architecture
itself provides.

The Ledger removes the failure mode at its root. There is exactly one canonical record
--- an ordered, append-only stream of atomic transactions --- and every other number
is computed on demand as a projection of that stream: balances, valuations,
profit-and-loss statements, balance sheets, regulatory reports. Two projections of one
record cannot disagree about any fact the record contains. That sentence is the
design's entire ambition; the rest of this specification is what it forces. Its scope
--- the recording, valuation, and lifecycle of trading-book positions held at fair
value, and the projection of committed activity into settlement --- is fixed
component by component in Chapter~17, together with the external authorities the
ledger reconciles against at its boundary and never replaces.

\subsection{The six commitments}

Six commitments are non-negotiable, and every design decision in this specification is
tested against them. Each is stated here with the one consequence it forces on the
rest of the document.

\textbf{Full auditability.} Every number is traceable, mechanically and completely, to
the recorded events that produced it. History is never rewritten: a mistake is
repaired by a recorded correction that names what it supersedes, so the error and its
repair both survive. This commitment forces the immutable, hash-chained log of
Chapter~\ref{ch:picture} and Chapter~3, and the repair discipline of Chapter~15: a
wrong booking is corrected by a compensating transaction through the same admission
door, in the open, never by amendment of the past.

\textbf{Full reproducibility.} Given the same record, any party --- the firm, a
counterparty, an auditor, a regulator --- recomputes the same state, the same views,
the same contract outputs, to the last minor unit. This commitment forces every arrow
of Chapter~\ref{ch:picture} to be a deterministic function of recorded inputs and
every quantity in this document to be an exact integer in minor units; the single
input not on the record, the clock, is confined to one machine, whose emissions are
themselves recorded (Chapter~4).

\textbf{Time travel embedded in the design.} The exact state of the system at any past
moment is reconstructible from the record alone, at any time, forever ---
independently of whether the software that produced the events still exists. This
commitment forces replay to re-read admitted transactions and never to re-execute
smart contracts; Chapter~14 states time travel as a theorem of the fold, not a feature
bolted onto it.

\textbf{Correctness.} Illegal states are made unrepresentable wherever possible: a
check can be forgotten or bypassed; a type cannot. Where wrongness cannot be prevented
--- a contract can always compute the wrong economics --- it is made mechanically
detectable and deterministically repairable. Defaults fail closed: the system stops
rather than improvises. This commitment forces the single admission door of
Chapter~4, the typed lifecycle graphs of Chapter~3, and the two layers of correctness
of Chapter~15 --- structural correctness guaranteed at admission, economic correctness
made checkable by recomputation.

\textbf{Testability.} The ledger is correct because it is built to be tested. Every
step is a pure, versioned function with explicit inputs, explicit outputs, and strong
types --- typing alone eliminates whole classes of defects before a single test runs
--- and every core invariant is expressed as an executable check over generated
products, events, and histories, never defended only in prose. Correctness then
composes: because the steps are pure and their composition is fixed --- the map, then
the fold --- verifying each part verifies the whole. This commitment forces the
invariant catalogue of Chapter~14 and the executable-check regime of Chapter~15, and
it forces a rule of method: a property that cannot be tested mechanically is
redesigned until it can be.

\textbf{Clarity.} One name per component and one component per name; behaviour carried
as declared data rather than hidden in code; every claim stated so that it can be
checked by a reader with no prior exposure. This commitment fixes the vocabulary of
this document --- unit, wallet, balance, move, transaction, watch, the immutable log,
projection, the Event Monitor, the Events Executor, the Transaction Executor, smart
contract, the market data operator, the three homes, virtual wallet, virtual ledger
--- and forces the declared-data discipline that recurs in every chapter: terms,
watches, market data operators, and collateral eligibility are data the record
carries, never behaviour a reader must infer from code.

The commitments are ordered by what arbitrates a conflict: a design that is auditable
and reproducible but untestable is redesigned; a design that is elegant but cannot be
replayed is rejected. No design decision in this document rests on custom or typical
practice; each is derived from the commitments and then, where an external standard
binds, shown to satisfy it. One boundary bounds all six: they secure one consistent
internal record, not its correspondence with the world outside --- internal
unbreakability is not external correctness, and whether a projection matches the market,
the custodian, or the counterparty is reconciled at the boundary
(Chapter~\ref{ch:marketdata}), never assumed.

\subsection{The discipline of the document}

The direction of authority is one-way. This specification is rated against the Ledger
Framework Constitution v1.1 and against nothing else; each chapter opens by naming the
constitution sections it discharges, and a genuine conflict with the constitution is
never drafted around --- it is recorded in Chapter~17's open-problems index with the
exact amendment it would require.

Six instruments run through the document as threads, introduced as a cast in
Chapter~\ref{ch:picture}: a cash-settled future, a one-touch option, a cash dividend,
a stock split, an OTC variance swap, and a total return swap. Their terms, parties,
dates, and quantities are fixed once and never vary, so a number met in one chapter
can be checked against the same number in another. Each visit is an \emph{episode} in
a fixed five-part shape: the trigger; the contract that catches it; the transactions
it produces, as numbered rows with exact quantities; what the door checks; and the
design point it substantiates. All quantities are exact integers in minor units;
amounts stated in whole currency are exact in cents.

Two conventions govern notation. Category-theoretic language appears only as a
\emph{second telling}, and no knowledge of category theory is assumed of the reader:
every concept is first fully explained in plain terms with a concrete example, and a
marked box may then restate it categorically to remove residual ambiguity for the reader
who wants it --- nothing later depends on any such box, and deleting them all
leaves the specification complete. Code appears only as short typed signatures and
data declarations; the reference implementation lies outside this document entirely.
